


%% Lecture 02 — Basic Nuclear Properties, Mass Formula, and Liquid Drop Model

\section{Lecture 02 — Basic Nuclear Properties, Mass Formula, and Liquid Drop Model}
\label{sec:lec02:nuclear_properties}

\subsection*{Overview}
This lecture transitions from the qualitative challenges of modeling the strong nuclear force to empirical methodologies used to extract fundamental nuclear properties. We begin by examining nuclear size and density profiles via electron scattering and form factors. We then exploit the charge-symmetry of the strong force in "mirror nuclei" to isolate Coulomb energy contributions. The lecture formalizes the definition of nuclear masses, binding energy, and nucleon separation energies, followed by a review of modern experimental mass-measurement techniques (e.g., MR-ToF, Penning Traps). Finally, these empirical observations are synthesized into the Semi-Empirical Mass Formula (SEMF) based on the Liquid Drop Model, allowing us to predict nuclear stability and decay pathways along isobaric chains.

\subsection*{1. Nuclear Size, Charge Distribution, and Form Factors}
To determine the internal structure and size of a nucleus, we rely on the scattering of high-energy electrons. 
Electrons are ideal probes because they are fundamental point particles (not subject to the strong force) and interact exclusively via the well-understood electromagnetic and weak interactions. To resolve the nucleus, the electron's de Broglie wavelength $\lambda$ must be on the order of the nuclear size ($\sim \text{fm}$), requiring beam energies around $100 \text{ MeV}$. At excessive energies, inelastic scattering occurs (creating new particles), which complicates structure probing.

In the framework of scattering theory, the transition matrix element is proportional to the Fourier transform of the potential. The spatial distribution of the nuclear charge $\rho(\vec{r})$ modifies the point-like scattering amplitude by a geometric factor known as the \textbf{Form Factor}, $F(\vec{q})$.

\subsubsection*{Derivation of the Form Factor}
Let $\vec{k}$ and $\vec{k}'$ be the initial and final wavevectors of the electron. The momentum transfer is defined as $\vec{q} = \vec{k} - \vec{k}'$. 
The potential generated by the nuclear charge distribution $\rho(\vec{r})$ dictates the scattering amplitude. Normalizing the charge distribution such that $\int \rho(\vec{r}) d^3r = Ze$, the form factor is defined as:
\begin{align*}
F(\vec{q}) &= \frac{1}{Ze} \int \rho(\vec{r}) e^{i \vec{q} \cdot \vec{r}} d^3r
\end{align*}
The experimentally measured differential cross-section is the point-charge (Mott/Rutherford) cross-section modulated by $|F(\vec{q})|^2$:
\begin{equation}
    \left( \frac{d\sigma}{d\Omega} \right)_{\text{exp}} = \left( \frac{d\sigma}{d\Omega} \right)_{\text{point}} \left| F(\vec{q}) \right|^2
    \label{eq:lec02:cross_section}
\end{equation}
By measuring the cross-section at various angles (and thus various $\vec{q}$), we can invert the Fourier transform to deduce $\rho(\vec{r})$. 
Experimental results reveal two universal constants for most nuclei:
1.  **Nuclear Density:** The core density is roughly constant across all nuclei:
    \begin{equation}\boxed{
        \rho_0 \approx 0.15 \text{ nucleons/fm}^3 \approx 2.5 \times 10^{14} \text{ g/cm}^3
        \label{eq:lec02:nuclear_density}
    }\end{equation}
2.  **Nuclear Radius:** Because the density is constant, the volume scales linearly with the mass number $A$, meaning the radius scales as:
    \begin{equation}\boxed{
        R = R_0 A^{1/3}
        \label{eq:lec02:nuclear_radius}
    }\end{equation}
    where $R_0 \approx 1.25 \text{ fm}$.

\subsection*{2. Mirror Nuclei and the Coulomb Energy Difference}
Mirror nuclei are pairs of isobars (same $A$) where the number of protons and neutrons are swapped, i.e., $^{A}_{Z}\text{X}_{N}$ and $^{A}_{Z-1}\text{Y}_{N+1}$ where $A = 2Z - 1$.
Because the strong nuclear force is approximately charge-symmetric (a proton-proton strong interaction is nearly identical to a neutron-neutron strong interaction), the primary difference in binding energy between mirror nuclei arises purely from the electrostatic Coulomb repulsion.

\subsubsection*{Derivation of Coulomb Energy Difference}
Assuming the nucleus is a uniformly charged sphere of radius $R$, the classical Coulomb self-energy is:
\begin{align*}\label{Eq:Coluomb Energy Difference}
E_C(Z) &= \frac{3}{5} \frac{(Z e)^2}{4\pi\epsilon_0 R}
\end{align*}
The energy difference $\Delta E_C$ between the mirror nuclei (going from $Z$ to $Z-1$) is:
\begin{align*}
\Delta E_C &= E_C(Z) - E_C(Z-1) \\
&= \frac{3}{5} \frac{e^2}{4\pi\epsilon_0 R} \left[ Z^2 - (Z-1)^2 \right] \\
&= \frac{3}{5} \frac{e^2}{4\pi\epsilon_0 R} (2Z - 1)
\end{align*}
Since $A = 2Z - 1$ for these specific mirror pairs, and substituting $R = R_0 A^{1/3}$, we get:
\begin{equation}
    \Delta E_C = \frac{3}{5} \frac{e^2}{4\pi\epsilon_0 R_0} A^{2/3}
    \label{eq:lec02:coulomb_diff}
\end{equation}
By measuring the mass (energy) difference between mirror nuclei, we can extract an independent empirical measurement of $R_0$.

\subsection*{3. Masses, Binding Energies, and Separation Energies}
In tables, tabulated masses are always **atomic masses** ($M_{\text{atom}}$), not bare nuclear masses ($M_{\text{nuc}}$). When calculating energetics, we must account for the electron masses.
\begin{align*}
M_{\text{nuc}}(A,Z) &\approx M_{\text{atom}}(A,Z) - Z m_e
\end{align*}
*(Note: We typically neglect the atomic electron binding energies, which are on the order of eV/keV, compared to nuclear energies on the order of MeV).*

\subsubsection*{Nuclear Binding Energy ($B$)}
The Binding Energy is the energy required to completely disassemble a nucleus into its constituent protons and neutrons:
\begin{align*}
B(A,Z) &= \left[ Z m_p + N m_n - M_{\text{nuc}}(A,Z) \right] c^2 \\
&= \left[ Z m_p + N m_n - (M_{\text{atom}}(A,Z) - Z m_e) \right] c^2 \\
\end{align*}
Grouping $m_p + m_e \approx m_H$ (the mass of a hydrogen atom), we write the highly practical formula:
\begin{equation}\boxed{
    B(A,Z) = \left[ Z m_H + N m_n - M_{\text{atom}}(A,Z) \right] c^2
    \label{eq:lec02:binding_energy}
}\end{equation}

\subsubsection*{Separation Energies and $Q$-values}
The **$Q$-value** dictates the energy released or absorbed in a reaction. For $X \rightarrow Y + c$, $Q = (M_X - M_Y - M_c)c^2$. Exothermic reactions have $Q > 0$.

The **Neutron Separation Energy ($S_n$)** is the energy required to remove one neutron (analogous to atomic ionization energy). It is the negative of the $Q$-value for neutron emission:
\begin{align*}
S_n &= \left[ M_{\text{atom}}(A-1, Z) + m_n - M_{\text{atom}}(A, Z) \right] c^2
\end{align*}
Similarly, the **Proton Separation Energy ($S_p$)** is:
\begin{align*}
S_p &= \left[ M_{\text{atom}}(A-1, Z-1) + m_H - M_{\text{atom}}(A, Z) \right] c^2
\end{align*}

\subsection*{4. Experimental Mass Measurement}
Given that typical binding energies per nucleon are $\sim 8 \text{ MeV}$, and nuclear masses are on the scale of GeV, precision mass measurement is critical ($\sim \text{keV}$ accuracy).
\begin{itemize}
    \item \textbf{Standard Mass Spectrometer:} Ions enter a Wien filter (crossed $\vec{E}$ and $\vec{B}$ fields). Only particles with velocity $v = \mathcal{E}/B$ pass un-deflected. They then enter a uniform magnetic field $\vec{B}'$, executing uniform circular motion where the Lorentz force provides the centripetal acceleration:
    \begin{align*}
    q v B' &= \frac{m v^2}{R} \implies R = \frac{m v}{q B'} \implies \frac{m}{q} = \frac{R B B'}{\mathcal{E}}
    \end{align*}
    By measuring the radius $R$, the mass-to-charge ratio is determined.
    \item \textbf{MR-ToF (Multi-Reflection Time-of-Flight):} Ions are trapped in an electrostatic potential well, bouncing back and forth. Because kinetic energy is uniform, velocity varies inversely with $\sqrt{m}$. Over many reflections, the flight time differences amplify: $t \propto \sqrt{m}$. This yields ultra-fast ($\sim$ ms), high-precision separation, ideal for short-lived exotic nuclei.
    \item \textbf{Penning Traps:} Ions are confined by a combination of a strong axial magnetic field and a quadrupole electric field. The precise measurement of the cyclotron frequency $\omega_c = qB/m$ via phase shifts allows for mass measurement accuracies down to $\Delta m / m \sim 10^{-10}$.
\end{itemize}

\subsection*{5. The Semi-Empirical Mass Formula (SEMF)}
To model the systematic trends in the nuclear binding energy (e.g., the saturation of $B/A \approx 8 \text{ MeV}$), we employ the **Liquid Drop Model**, treating the nucleus as an incompressible fluid droplet with several distinct physical contributions.

\begin{equation}\boxed{
    B(A,Z) = a_v A - a_s A^{2/3} - a_c \frac{Z(Z-1)}{A^{1/3}} - a_a \frac{(A-2Z)^2}{A} + \delta(A,Z)
    \label{eq:lec02:semf}
}\end{equation}

\textbf{Physical meaning of the terms:}
\begin{enumerate}
    \item \textbf{Volume Term ($+a_v A$):} Represents the strong nuclear force. Since $B/A$ is roughly constant, each nucleon only interacts with its nearest neighbors, meaning the binding energy scales linearly with the volume (and thus $A$).
    \item \textbf{Surface Term ($-a_s A^{2/3}$):} Nucleons on the surface have fewer neighbors than those in the bulk. This reduces the total binding energy proportional to the surface area $\propto R^2 \propto A^{2/3}$.
    \item \textbf{Coulomb Term ($-a_c Z(Z-1) A^{-1/3}$):} Electrostatic repulsion between the $Z$ protons. Scales as $1/R \propto A^{-1/3}$.
    \item \textbf{Asymmetry Term ($-a_a (A-2Z)^2 / A$):} Quantum mechanical in origin (Pauli Exclusion Principle). Nuclei favor $N = Z$. Any deviation costs energy.
    \item \textbf{Pairing Term ($\delta$):} Empirical observation that nucleons favor pairing up (spin up with spin down).
    \begin{align*}
        \delta(A,Z) = 
        \begin{cases} 
        +a_p A^{-1/2} & \text{even } Z, \text{ even } N \text{ (most stable)} \\
        0 & \text{odd } A \\
        -a_p A^{-1/2} & \text{odd } Z, \text{ odd } N \text{ (least stable)}
        \end{cases}
    \end{align*}
\end{enumerate}

\textbf{Example Calculation for $^{56}_{26}\text{Fe}$:}
Using the typical coefficients: $a_v = 15.5$, $a_s = 16.8$, $a_c = 0.72$, $a_a = 23$, $a_p = 12$ (all in MeV). For $^{56}\text{Fe}$, $Z=26, N=30$ (even-even $\implies$ positive pairing).
\begin{align*}
B(56, 26) &= 15.5(56) - 16.8(56)^{2/3} - 0.72\frac{26(25)}{56^{1/3}} - 23\frac{(56-52)^2}{56} + 12(56)^{-1/2} \\
&\approx 496 \text{ MeV}
\end{align*}
Thus, $B/A \approx \frac{496}{56} \approx 8.9 \text{ MeV}$, residing perfectly at the peak of the binding energy curve.

\subsubsection*{Isobaric Chains and Parabolas}
If we hold $A$ constant (an isobaric chain) and plot Mass vs. $Z$, the SEMF reveals a quadratic relationship: $M(A,Z) \propto \alpha Z^2 + \beta Z + \gamma$.
*   For **odd $A$**, the pairing term is zero. The plot is a single parabola. The minimum represents the single stable isobar. Nuclei higher on the walls undergo $\beta^-$ or $\beta^+$ decay to cascade down to the minimum.
*   For **even $A$**, the pairing term splits the plot into two nested parabolas (one for even-even, one for odd-odd). This can lead to multiple stable isobars for a given $A$ (as seen in $^{128}\text{Te}$ and $^{128}\text{Xe}$), where double-$\beta$ decay is theoretically possible but single $\beta$ decay is blocked by energy conservation.

\subsection*{Connections}
*   \textbf{Quantum Mechanics (Uncertainty Principle \& Atomic Size):} The radius of an atom is defined by the electron's properties. In the hypothetical case of Muonic Hydrogen (an electron swapped for a muon, $m_\mu \approx 207 m_e$), the Bohr radius shrinks by a factor of 207, drastically increasing the overlap of the wavefunction with the nucleus. This leads to the famous "Proton Radius Puzzle," where scattering data and muonic transition data yielded a $7\sigma$ discrepancy.
*   \textbf{Fourier Transforms:} The derivation of the form factor $F(\vec{q})$ is a direct application of the Born Approximation from QM scattering theory, analogous to X-ray diffraction in solid-state physics.

\subsection*{Exam / HW Hints}
*   \textbf{Always use Atomic Masses:} When calculating Binding Energy or $Q$-values, always use $m_H$ (Hydrogen mass) instead of $m_p$ (proton mass) if your table provides atomic masses. This automatically accounts for the $Z$ electron masses to first order.
*   \textbf{SEMF $Z(Z-1)$ vs $Z^2$:} Some textbooks approximate the Coulomb term with $Z^2$ instead of $Z(Z-1)$. $Z(Z-1)$ is strictly correct because a proton does not repel itself, but for large $Z$, the difference is negligible. Check which convention your professor prefers for numerical grading!